\subsection{Setup}
\label{subsec:exp_setup}





\paragraph{Datasets}



\begin{table}[!htbp]
\centering
\small
% \setlength{\tabcolsep}{4pt}
% \renewcommand{\arraystretch}{1}
\begin{tabular}{lcccc}
\toprule
\textbf{Dataset} & \textbf{Nodes} & \textbf{Edges} & \textbf{Max Degree} & \textbf{Ave. Degree} \\
\midrule
\textbf{AstroPh}   & 18,772  & 198,110  & 504   & 21.1 \\
\textbf{Enron}     & 36,692  & 138,831  & 1,383 & 7.6  \\
\textbf{Epinions1} & 75,879  & 405,740  & 3,044 & 10.7 \\
\textbf{Slashdot}  & 77,360  & 546,487  & 2,541 & 14.1 \\
\textbf{Twitter}   & 81,306  & 1,342,310 & 3,383 & 33.0 \\
\textbf{Epinions2} & 131,828 & 711,783  & 3,558 & 10.8 \\
\bottomrule
\end{tabular}
\caption{Summary of datasets used in the experiments.}
\label{tab:dataset}
\end{table}


We evaluate our mechanisms on six real-world graph datasets from SNAP~\cite{leskovec2014snap}, summarized in Table~\ref{tab:dataset}. 
\textbf{AstroPh} is a co-authorship network in astrophysics, 
\textbf{Enron} is an email communication network, 
\textbf{Twitter} and \textbf{Slashdot} are social networks derived from their respective platforms, and 
\textbf{Epinions1} and \textbf{Epinions2} represent a trust network and a social network from Epinions. 
All graphs were preprocessed to be undirected and free of self-loops.





\paragraph{Queries}


For $k$-line walk and path counting tasks, we evaluate cases with $k \in \{4, 5, {6}\}$.
The corresponding walk counting queries are denoted as $Q_{\mathrm{4\text{-}wk}}$, $Q_{\mathrm{5\text{-}wk}}$, {and $Q_{\mathrm{6\text{-}wk}}$},
and the path counting queries as $Q_{\mathrm{4\text{-}pt}}$, $Q_{\mathrm{5\text{-}pt}}$,
{and $Q_{\mathrm{6\text{-}pt}}$}.
For $k$-edge acyclic pattern counting, we consider {three} representative patterns, denoted $\pi_1$, $\pi_2$, {and $\pi_3$}, with edge counts $k = 4$, $k = 5$, {and $k=6$}, respectively. 
The associated queries are denoted by $Q_{\pi_1}$, $Q_{\pi_2}$, {and $Q_{\pi_3}$}.
For the first two patterns, we explore two distinct tree structures as pre-processing results: $\mathcal{T}_1$ and $\mathcal{T}_1'$ for $\pi_1$, and $\mathcal{T}_2$ and $\mathcal{T}_2'$ for $\pi_2$.
{For the third pattern, we consider a single tree structure $\mathcal{T}_3$.}
The structural details of the patterns and their corresponding trees are illustrated in Figure~\ref{fig:tree_form}.
Furthermore, we evaluated {three} cases of $k$-star counting with $k \in \{3, 4, {5}\}$, denoted as queries $Q_{3\star}$, $Q_{4\star}$, {and $Q_{5\star}$}.
All queries are evaluated without redundant counts in the results.


\begin{figure}[h]
  \centering
  \includegraphics[width=0.65\linewidth]{figures/exp_pattern.pdf} % file myfig.svg in figs/
  \caption{
  {Selected input patterns and their tree formulations with vertex orderings.}}
  \label{fig:tree_form}
\end{figure}




\paragraph{Experimental parameters}
All experiments are conducted on a Linux server with dual 64-core AMD EPYC 9555 CPUs and 2~TB RAM.
Each experiment is repeated $10$ times.
To evaluate utility, we discard the highest and lowest $2$ results and report the average of the remaining $6$ runs.
Communication cost is reported as the average over all $10$ runs.
We set the default privacy budget to $\varepsilon = 1$.
{All evaluated methods, including our parameter-free mechanisms, take only the privacy budget $\varepsilon$ and the query pattern as inputs, with no additional hyperparameters or tuning. 
For each query, all methods are run under identical settings to ensure fair comparison.}
{As $d(G)$ is assumed to be unknown in the experiments, both $\mathrm{WalkClip}_k$ and $\mathrm{LocalLap}_{k\star}$ require an additional round to estimate the noisy maximum degree $\tilde{d}(G)$. 
For these two methods, we allocate a privacy budget of $\varepsilon / 10$ to the estimation step, and the remaining $9\varepsilon / 10$ to the main counting algorithm. 
Following~\cite{betzer2024publishing}, we set the clipping factor\footnote{{The clipping threshold at round $\ell$ is defined as this factor raised to the power of $\ell$.}} in $\mathrm{WalkClip}_k$ as $\sqrt{\tilde{d}(G)}$.
}





\begin{table}[hbtp!]
\centering
% \small
% \setlength{\tabcolsep}{6pt}
% \renewcommand{\arraystretch}{0.96}
\begin{threeparttable}
\resizebox{0.95\columnwidth}{!}{%
\begin{tabular}{clccccccc}
\toprule
\textbf{Query} & \textbf{Method} & \textbf{Metric} &
\textbf{AstroPh} & \textbf{Enron} & \textbf{Epinions1} &
\textbf{Slashdot} & \textbf{Twitter} & \textbf{Epinions2} \\
\midrule
% ===================== k-line walks =====================
\multirow[c]{6}{*}[-2pt]{$Q_{\mathrm{4\text{-}wk}}$}
& \multirow[c]{2}{*}{$\mathcal{M}_{\mathrm{4\text{-}wk}}$}
& Rel. err (\%) & \cellcolor{gray!35}1.48 
& \cellcolor{gray!35}1.82
& \cellcolor{gray!35}1.18 
& \cellcolor{gray!35}1.75 
& \cellcolor{gray!35}1.22 
& \cellcolor{gray!35}0.91 \\
& & Comm. & \cellcolor{gray!12}6.76~MB
& \cellcolor{gray!12}7.01~MB 
& \cellcolor{gray!12}15.28~MB 
& \cellcolor{gray!12}17.27~MB 
& \cellcolor{gray!12}44.07~MB 
& \cellcolor{gray!12}26.72~MB \\

\cmidrule(lr){2-9}
& {\multirow[c]{2}{*}{$\mathrm{WalkClip}_4$}}
& {Rel. err (\%)} 
& {98.97} 
& {98.68} 
& {98.70} 
& {97.17} 
& {97.92} 
& {99.34} \\
& & {Comm.} & {10.50~GB} & {40.13~GB} & {171.60~GB} & {178.36~GB} & {197.02~GB} & {515.98~GB} \\

\cmidrule(lr){2-9}
& \multirow[c]{2}{*}{$RR_{\text{enum}}$}
& Rel. err (\%) & 572.86 & 929.78 & 733.06 & 1719.88 & 308.38 & 687.89 \\
& & Comm.$^\dagger$ & 168.00~MB & 641.95~MB & 2745.41~MB & 2853.63~MB & 3152.17~MB & 8255.56~MB \\
\midrule

\multirow[c]{5}{*}[-3pt]{$Q_{\mathrm{5\text{-}wk}}$}
& \multirow[c]{2}{*}{$\mathcal{M}_{\mathrm{5\text{-}wk}}$}
& Rel. err (\%) & \cellcolor{gray!35}2.27 
& \cellcolor{gray!35}2.30 
& \cellcolor{gray!35}1.72 
& \cellcolor{gray!35}2.80 
& \cellcolor{gray!35}1.49 
& \cellcolor{gray!35}0.88 \\
& & Comm. & \cellcolor{gray!12}10.07~MB 
& \cellcolor{gray!12}10.37~MB 
& \cellcolor{gray!12}22.63~MB 
& \cellcolor{gray!12}25.61~MB 
& \cellcolor{gray!12}65.79~MB 
& \cellcolor{gray!12}39.58~MB \\
\cmidrule(lr){2-9}

& {\multirow[c]{2}{*}{$\mathrm{WalkClip}_5$}}
& {Rel. err (\%)} 
& {99.76} 
& {99.59} 
& {99.61} 
& {98.88} 
& {99.35} 
& {99.84} \\
& & {Comm.} & {13.13~GB} & {50.16~GB} & {214.49~GB} & {222.95~GB} & {246.27~GB} & {644.98~GB} \\
\cmidrule(lr){2-9}


& \multirow[c]{1}{*}{$RR_{\text{enum}}$}
& Rel. err (\%) & 1854.25 & 3269.04 & 2346.32 & 8122.89 & 1027.45 & 2208.09 \\
\midrule

\multirow[c]{5}{*}[-4pt]{
% \marginnote{{R1.O4}}
{$Q_{\mathrm{6\text{-}wk}}$}}
& \multirow[c]{2}{*}{{$\mathcal{M}_{\mathrm{6\text{-}wk}}$}}
& {Rel. err (\%)} 
& \cellcolor{gray!35} {5.12} 
& \cellcolor{gray!35} {7.15} 
& \cellcolor{gray!35} {2.22} 
& \cellcolor{gray!35} {5.37} 
& \cellcolor{gray!35} {2.54} 
& \cellcolor{gray!35} {1.63} \\
& & {Comm.} 
& \cellcolor{gray!12} {13.38~MB} 
& \cellcolor{gray!12} {13.74~MB} 
& \cellcolor{gray!12} {29.97~MB} 
& \cellcolor{gray!12} {33.95~MB} 
& \cellcolor{gray!12} {87.51~MB} 
& \cellcolor{gray!12} {52.44~MB} \\
\cmidrule(lr){2-9}

& {\multirow[c]{2}{*}{$\mathrm{WalkClip}_6$}}
& {Rel. err (\%)} 
& {99.95} 
& {99.88} 
& {99.89} 
& {99.58} 
& {99.81} 
& {99.96} \\
& & {Comm.} & {15.75~GB} & {60.19~GB} & {257.39~GB} & {267.54~GB} & {295.53~GB} & {773.98~GB} \\


\cmidrule(lr){2-9}
& \multirow[c]{1}{*}{{$RR_{\text{enum}}$}}
& {Rel. err (\%)} 
& {6870.55} & {13021.01} & {8660.78} & {42530.83} & {3771.45} & {8179.33} \\
\midrule


% ===================== k-line paths =====================
\multirow[c]{3}{*}{$Q_{\mathrm{4\text{-}pt}}$}
& \multirow[c]{2}{*}{${\mathcal{M}}_{\mathrm{4\text{-}pt}}$}
& Rel. err (\%) & \cellcolor{gray!35}5.67
& \cellcolor{gray!35}11.47 
& \cellcolor{gray!35}5.95 
& \cellcolor{gray!35}14.02 
& \cellcolor{gray!35}8.91 
& \cellcolor{gray!35}8.79 \\
& & Comm. & \cellcolor{gray!12}3.55~MB 
& \cellcolor{gray!12}3.59~MB 
& \cellcolor{gray!12}7.84~MB
& \cellcolor{gray!12}8.92~MB 
& \cellcolor{gray!12}23.35~MB 
& \cellcolor{gray!12}13.76~MB \\
\cmidrule(lr){2-9}
& \multirow[c]{1}{*}{$RR_{\text{enum}}$}
& Rel. err (\%) & 477.46 & 791.34 & 600.68 & 1448.47 & 254.75 & 558.30 \\
\midrule

\multirow[c]{3}{*}{$Q_{\mathrm{5\text{-}pt}}$}
& \multirow[c]{2}{*}{${\mathcal{M}}_{\mathrm{5\text{-}pt}}$}
& Rel. err (\%) & \cellcolor{gray!35}11.77 
& \cellcolor{gray!35}5.95 
& \cellcolor{gray!35}11.55 
& \cellcolor{gray!35}10.23 
& \cellcolor{gray!35}4.13 
& \cellcolor{gray!35}9.55 \\
& & Comm. & \cellcolor{gray!12}3.58~MB 
& \cellcolor{gray!12}3.64~MB 
& \cellcolor{gray!12}7.96~MB 
& \cellcolor{gray!12}9.02~MB 
& \cellcolor{gray!12}23.54~MB
& \cellcolor{gray!12}13.92~MB \\
\cmidrule(lr){2-9}
& \multirow[c]{1}{*}{$RR_{\text{enum}}$}
& Rel. err (\%) & 1350.34 & 2400.53 & 1655.97 & 5877.49 & 733.03 & 1538.58 \\
\midrule

\multirow[c]{3}{*}{{$Q_{\mathrm{6\text{-}pt}}$}}
& \multirow[c]{2}{*}{{${\mathcal{M}}_{\mathrm{6\text{-}pt}}$}}
& {Rel. err (\%)} 
& \cellcolor{gray!35} {6.81} 
& \cellcolor{gray!35} {19.29} 
& \cellcolor{gray!35} {15.62} 
& \cellcolor{gray!35} {9.44} 
& \cellcolor{gray!35} {6.70} 
& \cellcolor{gray!35} {6.84} \\
& & {Comm.} 
& \cellcolor{gray!12} {3.60~MB} 
& \cellcolor{gray!12} {3.68~MB} 
& \cellcolor{gray!12} {8.02~MB} 
& \cellcolor{gray!12} {9.09~MB} 
& \cellcolor{gray!12} {23.59~MB}
& \cellcolor{gray!12} {14.03~MB} \\
\cmidrule(lr){2-9}
& \multirow[c]{1}{*}{{$RR_{\text{enum}}$}}
& {Rel. err (\%)} 
& {4115.66} & {7936.65} & {4977.07} & {25456.53} & {2209.38} & {4634.36} \\
\midrule


% ===================== acyclic patterns =====================
\multirow[c]{3}{*}{$Q_{\pi_1}$}
& \multirow[c]{2}{*}{$\mathcal{M}_{\mathcal{T}_1}$}
& Rel. err (\%) & \cellcolor{gray!35}12.10 & \cellcolor{gray!35}21.90 & \cellcolor{gray!35}11.71 & \cellcolor{gray!35}27.54 & \cellcolor{gray!35}15.16 & \cellcolor{gray!35}17.38 \\
& & Comm. & \cellcolor{gray!12}3.37~MB & \cellcolor{gray!12}3.37~MB & \cellcolor{gray!12}7.36~MB & \cellcolor{gray!12}8.39~MB & \cellcolor{gray!12}22.29~MB & \cellcolor{gray!12}12.89~MB \\
\cmidrule(lr){2-9}
& \multirow[c]{1}{*}{$RR_{\text{enum}}$}
& Rel. err (\%) & 256.42 & 198.42 & 178.08 & 368.72 & 65.25 & 166.18 \\
\midrule

\multirow[c]{3}{*}{$Q_{\pi_2}$}
& \multirow[c]{2}{*}{$\mathcal{M}_{\mathcal{T}_2}$}
& Rel. err (\%) & \cellcolor{gray!35}11.53 & \cellcolor{gray!35}23.27 & \cellcolor{gray!35}37.30 & \cellcolor{gray!35}48.14 & \cellcolor{gray!35}18.64 & \cellcolor{gray!35}17.24 \\
& & Comm. & \cellcolor{gray!12}3.32~MB & \cellcolor{gray!12}3.30~MB & \cellcolor{gray!12}7.23~MB & \cellcolor{gray!12}8.25~MB & \cellcolor{gray!12}21.96~MB & \cellcolor{gray!12}12.67~MB \\
\cmidrule(lr){2-9}
& \multirow[c]{1}{*}{$RR_{\text{enum}}$}
& Rel. err (\%) & 448.55 & 200.03 & 209.82 & 619.40 & 70.76 & 178.98 \\

\midrule 
\multirow[c]{3}{*}{{$Q_{\pi_3}$}}
& \multirow[c]{2}{*}{{$\mathcal{M}_{\mathcal{T}_3}$}}
& {Rel. err (\%)} 
& \cellcolor{gray!35} {18.74} 
& \cellcolor{gray!35} {48.23} 
& \cellcolor{gray!35} {23.81} 
& \cellcolor{gray!35} {77.08} 
& \cellcolor{gray!35} {28.20} 
& \cellcolor{gray!35} {43.45} \\
& & {Comm.} 
& \cellcolor{gray!12} {3.39~MB} 
& \cellcolor{gray!12} {3.40~MB} 
& \cellcolor{gray!12} {7.44~MB} 
& \cellcolor{gray!12} {8.46~MB} 
& \cellcolor{gray!12} {22.39~MB} 
& \cellcolor{gray!12} {13.00~MB} \\
\cmidrule(lr){2-9}
& \multirow[c]{1}{*}{{$RR_{\text{enum}}$}}
& {Rel. err (\%)} 
& {877.11}
& {273.41} 
& {205.49} 
& {547.55} 
& {54.35} 
& {232.20} \\

\bottomrule
\end{tabular}
}
\end{threeparttable}
\caption{
Relative error (\%) and communication cost (MB{/GB}) for $k$-line walk, $k$-line path, and $k$-edge acyclic pattern counting queries across six datasets.  
$^\dagger$Communication cost of $RR_{\mathrm{enum}}$ remains constant across queries.  
}
\label{tab:3queries}
\end{table}




\begin{table}[hbtp!]
\centering
% \small
% \setlength{\tabcolsep}{6pt}
% \renewcommand{\arraystretch}{0.9}
\resizebox{0.90\columnwidth}{!}{%
\begin{tabular}{clccccccc}
\toprule
\textbf{Query} & \textbf{Method} & \textbf{Metric} & \textbf{AstroPh} & \textbf{Enron} & \textbf{Epinions1} & \textbf{Slashdot} & \textbf{Twitter} & \textbf{Epinions2} \\
\midrule


\multirow[c]{4}{*}{$Q_{3\star}$}
& \multirow{2}{*}{$\mathcal{M}_{3\star}$} 
& Rel. err (\%) 
& \cellcolor{gray!35} 0.17 
& \cellcolor{gray!35} 0.09 
& \cellcolor{gray!35} 0.05 
& \cellcolor{gray!35} 0.06
& \cellcolor{gray!35} 0.04 
& \cellcolor{gray!35} 0.04 \\
& & Comm.$^{\dagger}$ 
& \cellcolor{gray!12} 0.14~MB 
& \cellcolor{gray!12} 0.28~MB 
& \cellcolor{gray!12} 0.58~MB
& \cellcolor{gray!12} 0.59~MB 
& \cellcolor{gray!12} 0.62~MB 
& \cellcolor{gray!12} 1.00~MB \\
\cmidrule(lr){2-9}


& \multirow{2}{*}{$\mathrm{LocalLap}_{3\star}$} & Rel. err (\%) & 6.79 & 8.16 & 19.52 & 16.47 & 5.64 & 10.62 \\
& & Comm.$^{\dagger}$ & 0.43~MB & 0.84~MB & 1.74~MB & 1.77~MB & 1.86~MB & 3.01~MB 


\\
\midrule

\multirow{2}{*}{$Q_{4\star}$}
& $\mathcal{M}_{4\star}$ & Rel. err (\%) & 
\cellcolor{gray!35}0.49 & 
\cellcolor{gray!35}0.19 & 
\cellcolor{gray!35}0.15 & 
\cellcolor{gray!35}0.06 & 
\cellcolor{gray!35}0.09 & 
\cellcolor{gray!35}0.09 \\
\cmidrule(lr){2-9}
& $\mathrm{LocalLap}_{4\star}$ & Rel. err (\%) & 
33.88 & 16.69 & 43.04 & 42.78 & 16.55 & 35.67 \\

\midrule

\multirow{2}{*}{{$Q_{5\star}$}}
& {$\mathcal{M}_{5\star}$} & {Rel. err (\%)} & \cellcolor{gray!35}
{0.50} & \cellcolor{gray!35}{0.22} & \cellcolor{gray!35}{0.18} & \cellcolor{gray!35}{0.14} & \cellcolor{gray!35}{0.11} & \cellcolor{gray!35}{0.09} \\
\cmidrule(lr){2-9}
& {$\mathrm{LocalLap}_{5\star}$} & {Rel. err (\%)} & 
{38.58} & {24.81} & {82.03} & {53.66} & {26.92} & {37.21} \\


\bottomrule
\end{tabular}
}
\caption{Relative error (\%) and communication cost (MB) for $k$-star queries on six datasets. $^{\dagger}$Communication cost of both methods remains constant across star counting queries.}
\label{tab:star}
\end{table}

